%=====================================================================
\chapter{Dynamic Programming II: Classic Problems}\label{ch:dptwo}
%=====================================================================
\locator{4}{Part IV --- The Design Paradigms}

\begin{outcomes}
By the end of this chapter you will be able to:
\begin{tight}
  \item \textbf{Derive} the recurrence for longest common subsequence and edit
        distance, and \textbf{implement} both \emph{(CLO-6)}.
  \item \textbf{Reduce} a dynamic program's space from $\bigTh{mn}$ to
        $\bigTh{\min(m,n)}$, and \textbf{state} what that costs.
  \item \textbf{Formulate} 0/1 knapsack as a dynamic program and
        \textbf{analyse} its $\bigTh{nW}$ running time.
  \item \textbf{Explain} why $\bigTh{nW}$ is \emph{pseudo-polynomial} rather
        than polynomial, in terms of input size.
  \item \textbf{Choose} a DP formulation by identifying the subproblem
        parameters.
\end{tight}
\end{outcomes}

\begin{prereq}
\chref{dpone} for the method. \chref{cases} is essential for Section~16.4 ---
the pseudo-polynomial argument is exactly that chapter's point about a number's
size being its bit length.
\end{prereq}

\begin{keyterms}
longest common subsequence $\bullet$ edit distance $\bullet$ Levenshtein
$\bullet$ alignment $\bullet$ 0/1 knapsack $\bullet$ matrix-chain
multiplication $\bullet$ pseudo-polynomial $\bullet$ rolling array $\bullet$
reconstruction $\bullet$ subproblem parameters
\end{keyterms}

\section{Longest Common Subsequence}

\begin{definitionbox}[Subsequence, LCS]
A \term{subsequence} of a sequence is obtained by deleting zero or more
elements without reordering the rest. \texttt{ALRIT} is a subsequence of
\texttt{ALGORITHM}.

\smallskip
The \term{longest common subsequence} of $X$ and $Y$ is a longest sequence that
is a subsequence of both. Note that it need not be contiguous --- that would be
the longest common \emph{substring}, a different and easier problem.
\end{definitionbox}

\begin{conceptbox}[Finding the recurrence]
Let $c[i][j]$ be the LCS length of the prefixes $X_{1..i}$ and $Y_{1..j}$.
Consider the last characters.

\smallskip
\textbf{If $x_{i} = y_{j}$}, there is an LCS ending with that character ---
including it can never hurt --- so $c[i][j] = c[i-1][j-1] + 1$.

\smallskip
\textbf{If $x_{i} \ne y_{j}$}, the LCS cannot end with both, so at least one of
$x_{i}, y_{j}$ is unused, and $c[i][j] = \max(c[i-1][j],\, c[i][j-1])$.

\smallskip
\[ c[i][j] = \begin{cases}
   0 & i = 0 \text{ or } j = 0, \\
   c[i-1][j-1] + 1 & x_{i} = y_{j}, \\
   \max\bigl(c[i-1][j],\, c[i][j-1]\bigr) & \text{otherwise.}
   \end{cases} \]

\smallskip
\textbf{Both conditions of \chref{dpone} hold.} Optimal substructure is the
argument above. Overlap is severe: $c[i-1][j]$ and $c[i][j-1]$ both depend on
$c[i-1][j-1]$, so a naive recursion solves it twice, and the duplication
compounds exponentially.
\end{conceptbox}

\dsfig{\aoaalg{\algname{LCS-Length}$(X, Y, m, n)$}{
\For{$i \gets 0$ \textbf{to} $m$}
  \State $c[i][0] \gets 0$
\EndFor
\For{$j \gets 0$ \textbf{to} $n$}
  \State $c[0][j] \gets 0$
\EndFor
\For{$i \gets 1$ \textbf{to} $m$}
  \For{$j \gets 1$ \textbf{to} $n$}
    \If{$x_{i} = y_{j}$}
      \State $c[i][j] \gets c[i-1][j-1] + 1$
    \Else
      \State $c[i][j] \gets \max\bigl(c[i-1][j],\, c[i][j-1]\bigr)$
    \EndIf
  \EndFor
\EndFor
\State \Return $c[m][n]$
}}{\algname{LCS-Length}. $mn$ subproblems, $\bigO{1}$ work each, so
$\bigTh{mn}$ time --- and $\bigTh{mn}$ space, which Section~16.3 reduces.}{lcs}

\begin{theorem}[LCS]\label{thm:lcs}
\algname{LCS-Length} computes the LCS length of $X$ and $Y$ in $\bigTh{mn}$
time and $\bigTh{mn}$ space.
\end{theorem}

\begin{proof}
The nested loops perform $mn$ iterations with $\bigO{1}$ bodies, and the
initialisation is $\bigTh{m+n}$, giving $\bigTh{mn}$ time; the table gives the
space.

\smallskip
Correctness follows by induction on $i+j$. For $i = 0$ or $j = 0$ one prefix is
empty and the LCS length is 0. For $i, j \ge 1$, assume the entries with
smaller $i+j$ are correct. If $x_{i} = y_{j}$: any common subsequence of
$X_{1..i}, Y_{1..j}$ not ending in that character can be extended by it, so
some LCS does end in it, and removing it leaves an LCS of $X_{1..i-1},
Y_{1..j-1}$ --- giving $c[i-1][j-1]+1$. If $x_{i} \ne y_{j}$: an LCS cannot use
both last characters, so it is an LCS of $X_{1..i-1}, Y_{1..j}$ or of
$X_{1..i}, Y_{1..j-1}$, and the maximum of the two is correct.
\end{proof}

\begin{worked}[A Small Instance, Checkable by Eye]
\texttt{code/ch16\_dp2.cpp} on \texttt{ALGORITHM} and \texttt{ALTRUISTIC}:
\begin{center}
LCS length $= 5$, the subsequence \texttt{ALRIT};\qquad
edit distance $= 6$.
\end{center}
Trace \texttt{ALRIT} through both words to confirm it is a subsequence of each,
and satisfy yourself that no common subsequence of length 6 exists. Small
instances you can check by hand are worth more than large ones you cannot.
\end{worked}

\section{Edit Distance}

\begin{definitionbox}[Edit Distance]
The \term{edit distance} (Levenshtein distance) between two strings is the
minimum number of single-character insertions, deletions and substitutions that
turn one into the other.
\end{definitionbox}

\begin{conceptbox}[The same table, three choices instead of two]
Let $d[i][j]$ be the distance between $X_{1..i}$ and $Y_{1..j}$. To account for
the last position there are exactly three possibilities:
\[ d[i][j] = \min\begin{cases}
   d[i-1][j] + 1 & \text{delete } x_{i}, \\
   d[i][j-1] + 1 & \text{insert } y_{j}, \\
   d[i-1][j-1] + [\,x_{i} \ne y_{j}\,] & \text{substitute, or match free.}
   \end{cases} \]
with $d[i][0] = i$ and $d[0][j] = j$ --- deleting or inserting everything.

\smallskip
\textbf{LCS and edit distance are the same algorithm with a different cost
table}, which is worth noticing: once you can write one, you can write
spell-checking, DNA alignment, \texttt{diff}, and fuzzy search. All of them are
$\bigTh{mn}$ with different weights on the three moves.
\end{conceptbox}

\begin{worked}[Measured: Both Are Quadratic]
Random strings over a 4-letter alphabet; doubling both lengths should quadruple
the time.

\rowcolors{2}{white}{lightbg}
\begin{center}
\begin{tabular}{@{}rrrrrr@{}}
\toprule
\rowcolor{primary!15}
$n$ & LCS full (ms) & ratio & LCS 2-row (ms) & ratio & edit dist (ms) \\
\midrule
   500 &   1.4 & ---  &   0.7 & ---  &   1.3 \\
1{,}000 &   5.5 & 4.08 &   3.2 & 4.65 &   5.2 \\
2{,}000 &  24.4 & 4.43 &  14.6 & 4.59 &  22.1 \\
4{,}000 & 107.0 & 4.39 &  61.2 & 4.19 &  82.4 \\
8{,}000 & 403.6 & 3.77 & 240.2 & 3.93 & 320.3 \\
\bottomrule
\end{tabular}
\end{center}

\smallskip
Mean ratios 4.17 and 4.34 against a predicted 4. The two-row version is
consistently about 1.7$\times$ faster than the full table --- not because it
does less arithmetic, but because it touches $\bigTh{n}$ memory instead of
$\bigTh{n^{2}}$ and stays in cache.
\end{worked}

\section{Saving Space, and What It Costs}

\begin{conceptbox}[The rolling array]
Row $i$ of the LCS table depends only on row $i-1$ and on earlier entries of
row $i$. So only two rows need exist at a time, and the space falls from
$\bigTh{mn}$ to $\bigTh{n}$ --- or $\bigTh{\min(m,n)}$, by iterating over the
longer string in the outer loop.

\smallskip
At $m = n = 100{,}000$ that is the difference between 40 GB and 400 KB. It is
not an optimisation; it is the difference between possible and impossible.

\smallskip
\textbf{And here is the cost.} Reconstructing the actual subsequence requires
walking \emph{backwards} through the table, which needs the whole table. With
two rows you get the length and nothing else.

\smallskip
\textbf{So the choice is:}
\begin{tight}
  \item need the length only $\Rightarrow$ two rows, $\bigTh{\min(m,n)}$ space;
  \item need the sequence $\Rightarrow$ full table, $\bigTh{mn}$ space;
  \item need the sequence and cannot afford the table $\Rightarrow$
        Hirschberg's algorithm, which combines the rolling array with
        divide-and-conquer to reconstruct in $\bigTh{mn}$ time and
        $\bigTh{\min(m,n)}$ space. It is the right answer and it is twice the
        code.
\end{tight}
\end{conceptbox}

\section{0/1 Knapsack, and a Word That Matters}

\begin{definitionbox}[0/1 Knapsack]
Given $n$ items, item $i$ having integer weight $w_{i}$ and value $v_{i}$, and
a knapsack of integer capacity $W$, choose a subset of maximum total value
whose total weight is at most $W$. Each item is taken whole or not at all ---
hence 0/1.
\end{definitionbox}

\dsfig{\aoaalg{\algname{Knapsack}$(w, v, n, W)$}{
\For{$\mathit{cap} \gets 0$ \textbf{to} $W$}
  \State $b[\mathit{cap}] \gets 0$
\EndFor
\For{$i \gets 1$ \textbf{to} $n$}
  \For{$\mathit{cap} \gets W$ \textbf{downto} $w_{i}$}
    \State $b[\mathit{cap}] \gets \max\bigl(b[\mathit{cap}],\;
           b[\mathit{cap} - w_{i}] + v_{i}\bigr)$
  \EndFor
\EndFor
\State \Return $b[W]$
}}{\algname{Knapsack}, one row at a time. The inner loop counts \emph{down},
which is what enforces 0/1: counting up would let item $i$ be used more than
once, solving the unbounded knapsack problem instead.}{knapsack}

\begin{theorem}[Knapsack]\label{thm:knapsack}
\algname{Knapsack} computes the optimal value in $\bigTh{nW}$ time and
$\bigTh{W}$ space.
\end{theorem}

\begin{proof}
There are $nW$ subproblems --- one per (item prefix, capacity) pair --- each
resolved by one comparison, giving $\bigTh{nW}$ by the formula of
\chref{dpone}. The single array gives $\bigTh{W}$ space.

\smallskip
For the downward inner loop: when $b[\mathit{cap}]$ is updated it reads
$b[\mathit{cap}-w_{i}]$, which at that moment still holds the value from
iteration $i-1$, because indices below $\mathit{cap}$ have not yet been visited
in this pass. So each item is considered at most once per capacity, as 0/1
requires.
\end{proof}

\begin{worked}[Checked Against an Oracle]
\chref{bruteforce}'s technique again: 300 random instances with $n = 18$, where
all $2^{18}$ subsets can be enumerated.
\begin{center}
\textbf{300 instances, 0 disagreements.}
\end{center}
Worth the twenty lines. The downward loop is exactly the kind of detail that is
easy to get backwards and produces plausible-but-wrong answers.
\end{worked}

\begin{alertbox}[$\bigTh{nW}$ is not a polynomial running time]
This is the most important sentence in the chapter, and it takes
\chref{cases} to understand.

\smallskip
$\bigTh{nW}$ \emph{looks} polynomial: no exponents, no factorials. But a
running time is polynomial when it is polynomial in the \textbf{size of the
input}, and \chref{cases} established that a number's size is its \emph{bit
length}, not its value.

\smallskip
$W$ is one number. Writing it takes $b \approx \lg W$ bits. So
\[ \bigTh{nW} = \bigTh{n \cdot 2^{b}} \]
--- \textbf{exponential in the length of the input}. Add one bit to $W$ and the
running time doubles.

\smallskip
An algorithm polynomial in the numeric value of its input but exponential in
its length is called \term{pseudo-polynomial}. \chref{complexity} explains why
this distinction decides whether a problem counts as solved: 0/1 knapsack is
NP-complete, and this algorithm is not a counter-example to that.
\end{alertbox}

\begin{worked}[The Claim, Measured]
$n$ fixed at 200. Only the capacity changes, and the capacity is a single
number.

\rowcolors{2}{white}{lightbg}
\begin{center}
\begin{tabular}{@{}rrrr@{}}
\toprule
\rowcolor{primary!15}
$W$ & bits & time (ms) & ratio \\
\midrule
      65{,}536 & 17 &    16.4 & ---  \\
     131{,}072 & 18 &    34.9 & 2.13 \\
     262{,}144 & 19 &    64.1 & 1.84 \\
     524{,}288 & 20 &   134.0 & 2.09 \\
   1{,}048{,}576 & 21 &   275.9 & 2.06 \\
   2{,}097{,}152 & 22 &   630.5 & 2.29 \\
   4{,}194{,}304 & 23 & 1{,}138.3 & 1.81 \\
   8{,}388{,}608 & 24 & 2{,}278.9 & 2.00 \\
  16{,}777{,}216 & 25 & 4{,}540.1 & 1.99 \\
\bottomrule
\end{tabular}
\end{center}

\smallskip
\textbf{Eight ratios, mean 2.02. Each row adds exactly one bit of input and
doubles the running time.}

\smallskip
\textbf{Compare this table with \chref{cases}'s trial-division table}, where
adding two bits to a prime doubled the work, mean ratio 2.02. They are the same
phenomenon, reached from opposite directions: there, an algorithm that looked
exponential and was; here, an algorithm that looks polynomial and is not.

\smallskip
\textbf{Extrapolate.} A 64-bit capacity --- an entirely ordinary thing to find
in a struct --- is $2^{39}$ times the last row: about $7 \times 10^{11}$
seconds, or \textbf{22{,}000 years}. The algorithm is excellent for small
capacities and useless for large ones, and no amount of optimisation changes
that, because the table has $nW$ entries and they must be written.
\end{worked}

\begin{conceptbox}[How to find the subproblem parameters]
Every DP in this chapter is defined by \emph{what the subproblems are indexed
by}, and finding that is the whole design task.

\rowcolors{2}{white}{lightbg}
\begin{longtable}{@{}L{3.4cm}L{4.2cm}L{3.0cm}L{3.4cm}@{}}
\toprule
\rowcolor{primary!15}
\textbf{Problem} & \textbf{Subproblem indexed by} & \textbf{Count} &
\textbf{Time} \\
\midrule
\endhead
rod cutting & remaining length & $n$ & $\bigTh{n^{2}}$ \\
LCS & two prefix lengths & $mn$ & $\bigTh{mn}$ \\
edit distance & two prefix lengths & $mn$ & $\bigTh{mn}$ \\
0/1 knapsack & item prefix, capacity & $nW$ & $\bigTh{nW}$ \\
matrix chain & a range $(i,j)$ & $n^{2}$ & $\bigTh{n^{3}}$ \\
\bottomrule
\end{longtable}

\smallskip
Matrix-chain multiplication is the one with $\bigTh{n^{3}}$: there are
$\bigTh{n^{2}}$ ranges and each requires a maximum over $\bigO{n}$ split
points. The formula of \chref{dpone} gives the time directly once the
parameters are named.
\end{conceptbox}

\begin{pitfall}
\begin{tight}
  \item \textbf{Calling $\bigTh{nW}$ polynomial.} It is pseudo-polynomial, and
        measured, one more bit of $W$ doubles the time.
  \item \textbf{Counting the knapsack inner loop upwards.} That solves
        \emph{unbounded} knapsack --- each item reusable. Both are legitimate
        problems; they are not the same problem.
  \item \textbf{Confusing subsequence with substring.} LCS allows gaps; longest
        common substring does not and has a different recurrence.
  \item \textbf{Using the two-row trick and then trying to reconstruct.} The
        table is gone. Use the full table or Hirschberg's algorithm.
  \item \textbf{Getting the edit-distance base cases wrong.} $d[i][0] = i$,
        not 0 --- deleting $i$ characters costs $i$.
  \item \textbf{Skipping the brute-force oracle.} 300 instances at $n = 18$
        cost seconds and caught nothing here, which is the point: it is cheap
        evidence that the loop direction is right.
  \item \textbf{Assuming a DP's space equals its time.} LCS is $\bigTh{mn}$
        time and can be $\bigTh{\min(m,n)}$ space.
\end{tight}
\end{pitfall}

\begin{lab}[Three Classics and One Surprise]
Reproduces all four tables. Expect thirty minutes; part D takes about nine
seconds in its last row.

\begin{lstlisting}[language=C++]
// ch16_dp2.cpp -- the classic dynamic programs, and one nasty surprise.
// Build: cl /O2 /EHsc /std:c++17 ch16_dp2.cpp        (see Appendix A)

// --- 1. LCS: the full table, so the subsequence can be recovered ---
    for (size_t i = 1; i <= m; ++i)
        for (size_t j = 1; j <= n; ++j)
            c[i][j] = (a[i-1] == b[j-1]) ? c[i-1][j-1] + 1
                                         : std::max(c[i-1][j], c[i][j-1]);

// --- 2. LCS in two rows --------------------------------------------
// Row i depends only on row i-1, so only two rows need to exist. The
// LENGTH survives; the subsequence does not, because reconstructing
// it needs the whole table. That is the trade, and it is typical.

// --- 3. Edit distance: the same shape, three choices per cell ------
// Insert, delete or substitute. The recurrence differs from LCS only
// in what the three predecessors cost.
            cur[j] = std::min({prev[j] + 1,            // delete
                               cur[j - 1] + 1,         // insert
                               prev[j-1] + (a[i-1] != b[j-1])});

// --- 4. 0/1 knapsack: Theta(nW), and W is a number not a count -----
// The inner loop counts DOWN, which is what enforces 0/1: counting up
// would let item i be used more than once.
static long long knapsack(const std::vector<int>& w,
                          const std::vector<long long>& v, int W) {
    std::vector<long long> best(W + 1, 0);
    for (size_t i = 0; i < w.size(); ++i)
        for (int cap = W; cap >= w[i]; --cap)
            best[cap] = std::max(best[cap], best[cap - w[i]] + v[i]);
    return best[W];
}

// --- 5. The same problem by brute force, as an oracle --------------
// 2^n subsets. Only usable for n <= 22, which is exactly why DP
// exists -- and exactly why it is a trustworthy oracle.

// --- 6. LCS and edit distance are Theta(mn) ------------------------
// --- 7. A worked instance, small enough to check by eye ------------
// --- 8. Knapsack: DP against brute force ---------------------------

// --- 9. The surprise: doubling W doubles the time ------------------
// n is FIXED at 200 here. Only the capacity changes, and the capacity
// is one number. Writing 2W instead of W adds ONE BIT to the input
// and doubles the running time. Theta(nW) is polynomial in the VALUE
// of W and exponential in its LENGTH: pseudo-polynomial.
\end{lstlisting}

\textbf{What to notice.} Part 9's \texttt{bits} column is the whole lesson. The
input is not getting appreciably bigger --- it gains one bit per row, while the
200 items stay fixed --- and the running time doubles each time.

\smallskip
\textbf{Then try to break it.} Change the inner loop of \texttt{knapsack} to
count \emph{up} and re-run part 8 against the brute-force oracle. It will
disagree, and the answers will be too \emph{large} --- the algorithm now takes
some items repeatedly. Then work out which problem the upward loop solves, and
note that it is a perfectly good algorithm for a different question.
\end{lab}

\begin{checkpoint}
\begin{tightnum}
  \item Write the LCS recurrence and explain why the case $x_{i} = y_{j}$ does
        not need a maximum.
  \item Explain why $\bigTh{nW}$ is pseudo-polynomial, using the measured
        doubling table, and relate it to \chref{cases}'s trial division.
  \item Give one problem for which the two-row space trick is appropriate and
        one for which it is not, and say what distinguishes them.
\end{tightnum}
\end{checkpoint}

\begin{chaptersummary}
\begin{tight}
  \item LCS and edit distance share one table shape: $mn$ subproblems indexed
        by two prefix lengths, $\bigO{1}$ work each, so $\bigTh{mn}$. They
        differ only in the cost of the three moves.
  \item Measured: both quadratic, mean ratios 4.17 and 4.34 on doubling against
        a predicted 4; the two-row version was 1.7$\times$ faster through cache
        behaviour alone.
  \item The rolling array cuts space from $\bigTh{mn}$ to $\bigTh{\min(m,n)}$
        and makes reconstruction impossible. Hirschberg's algorithm gets both,
        at twice the code.
  \item 0/1 knapsack is $\bigTh{nW}$ with $nW$ subproblems. The inner loop must
        count down, or it solves the unbounded problem instead.
  \item $\bigTh{nW}$ is \textbf{pseudo-polynomial}: $W$ is a value, its input
        size is $\lg W$ bits, so the time is $\bigTh{n2^{b}}$.
  \item Measured with $n$ fixed at 200: eight doublings of $W$ gave time ratios
        of mean 2.02, each row adding one bit of input. A 64-bit capacity would
        take about 22{,}000 years.
  \item Find the subproblem parameters first; the running time follows from
        (count) $\times$ (work each).
\end{tight}
\end{chaptersummary}

\begin{reviewq}
  \item \textbf{(MCQ)} The LCS of two strings of lengths $m$ and $n$ is computed
        in: (a)~$\bigTh{m+n}$; (b)~$\bigTh{mn}$; (c)~$\bigTh{mn\lg n}$;
        (d)~$\bigTh{2^{n}}$.
  \item \textbf{(MCQ)} $\bigTh{nW}$ for 0/1 knapsack is:
        (a)~polynomial; (b)~pseudo-polynomial; (c)~logarithmic;
        (d)~constant in $W$.
  \item \textbf{(Short)} Derive the edit-distance recurrence, naming the three
        operations and their costs \emph{(CLO-6)}.
  \item \textbf{(Short)} Why must the knapsack inner loop count downwards?
        State what the upward version computes.
  \item \textbf{(Short)} Explain the trade-off between the full LCS table and
        the two-row version.
  \item \textbf{(Applied)} Give the dynamic program for matrix-chain
        multiplication: state the subproblem parameters, the recurrence, the
        number of subproblems and the running time.
  \item \textbf{(Applied)} A colleague proposes shipping the $\bigTh{nW}$
        knapsack solver for a logistics system where capacities are given in
        grams and can reach $10^{8}$. Using the measured table, estimate the
        running time, and propose two alternatives.
\end{reviewq}
